% =====================================================================
% Section 2 — Background and Related Work
% =====================================================================
\section{Background and Related Work}\label{sec:related}

\subsection{The JVM web-framework landscape}

Server-side Java web development is dominated by a small number of
mature frameworks that occupy distinct points along two axes: how much
they bundle, and when they do their work (at runtime versus at
compile/build time).

\emph{Spring Boot}~\cite{springboot} is the reference
``batteries-included'' stack: extensive auto-configuration, a vast
starter ecosystem, and runtime dependency injection driven by
reflection and classpath scanning. Its productivity is paid for in
artifact size and start-up cost. \emph{Quarkus}~\cite{quarkus} and
\emph{Micronaut}~\cite{micronaut} were designed to reduce exactly those
costs by moving dependency injection and metadata processing to
\emph{compile time}, which shrinks reflection at runtime and makes both
frameworks first-class targets for GraalVM native-image~\cite{graalvm}
ahead-of-time compilation. \emph{Javalin}~\cite{javalin} takes the
opposite, minimalist route: a thin, explicit API over an embedded
Jetty server, with relatively few dependencies but delegation of most
capabilities (JSON, persistence) to external libraries chosen by the
developer.

What unites all four is that non-trivial capabilities---JSON
serialisation, connection pooling, scheduling, and above all
security---are provided by \emph{third-party libraries}, whether
bundled by default or added by the developer. Security in particular is
predominantly \emph{opt-in}: a separate module (e.g.\ a security
starter) that must be added and configured before protections apply.

\subsection{Micro-frameworks and zero-dependency designs}

A long line of ``micro'' frameworks (in the JVM world and beyond)
reduces footprint by shipping only a routing core. This lowers the
dependency count but rarely reaches \emph{zero}: JSON and pooling are
still delegated, and the security posture remains additive. JxMVC
differs by implementing these subsystems in-framework on top of the JDK
and Jakarta Servlet API, so that the runtime dependency count is zero
(excluding the JDBC driver) and the compiled core remains small enough
to audit in full.

\subsection{Security-by-default}

The principle of \emph{secure defaults}---systems should be safe out of
the box, with the burden on relaxing rather than adding
protections---is long established in security engineering and codified
in practitioner guidance such as the OWASP Top~10~\cite{owasp}. In
mainstream frameworks, however, defaults are typically permissive for
convenience, and hardening (security headers, CSRF tokens, rate
limiting, body-size caps, anti open-redirect) is assembled by the
developer. JxMVC's distinguishing design decision is to apply these
protections \emph{in the request pipeline itself}, so that they hold
for every request unless explicitly and locally relaxed.

\subsection{Positioning}

Table~\ref{tab:qualitative} summarises the design axes along which the
evaluated frameworks differ. JxMVC is not proposed as a faster or more
capable replacement for these systems in the general case; it is a
different \emph{point in the design space}---zero runtime dependencies,
in-framework subsystems, and security-by-default---aimed at deployments
where deployable size, auditability, and safe defaults outweigh peak
throughput or minimal memory.

% Qualitative comparison (Section 3). Defensible, high-level claims only.
% Requires: booktabs. Keep textual (no pifont) for portability.
\begin{table*}[t]
  \centering
  \caption{Qualitative comparison of the evaluated frameworks. ``Own''
    means the capability is implemented in-framework without a
    third-party runtime library; ``Lib'' means it is delegated to a
    bundled external library. Security-by-default denotes whether the
    request pipeline applies protections (headers, CSRF, rate limiting,
    body caps) without opt-in configuration.}
  \label{tab:qualitative}
  \small
  \begin{tabular}{@{}lccccc@{}}
    \toprule
    & \textbf{JxMVC} & \textbf{Spring Boot} & \textbf{Quarkus} & \textbf{Micronaut} & \textbf{Javalin} \\
    \midrule
    Runtime dependencies   & \textbf{0}        & many        & many        & moderate    & few \\
    JSON                   & Own               & Lib (Jackson)& Lib        & Lib (Jackson)& Lib (Jackson) \\
    Connection pool        & Own (CAS)         & Lib (Hikari)& Lib (Agroal)& Lib (Hikari)& external \\
    Dependency injection   & Own (reflection)  & runtime     & compile-time& compile-time& none \\
    Security-by-default    & \textbf{Yes (pipeline)} & opt-in & opt-in     & opt-in      & manual \\
    OIDC (JDK-only)        & \textbf{Yes}      & Lib         & Lib         & Lib         & manual \\
    GraalVM native         & not yet           & partial (AOT)& first-class& first-class & limited \\
    Core artifact size     & \textbf{253\,KB}  & MB-scale    & MB-scale    & MB-scale    & MB-scale \\
    \bottomrule
  \end{tabular}
\end{table*}

